% Project-specific derivation for the retained literal compiler.
% This fragment makes no claim of generic quantum lower-bound novelty.
\section{Excluding still wider addresses within the fixed compiler}
\label{sec:raw-wider-address-exclusion}

I distinguish an exclusion within a fixed emitted implementation from an
algorithmic complexity lower bound. In this section the classical financial
models, $q=28$, narrow exponential span $b_0=1/64$, both payoff graphs,
normal widths $\{8,10,12\}$, outer precisions, two coherent unary-lookup
architectures and canonical-QPE controller are fixed. The scalar address
width is denoted by $b_e$, rather than by $b_0$.

The dyadic input interval contains $2^{q-6}+1$ representable arguments.
An address width $b_e$ gives $N=2^{b_e}$ cells, with remainder width
$q-6-b_e=22-b_e$. Hence this interpolator requires $1\leq b_e\leq22$.
The widths $11,12,13,14$ leave respectively $11,10,9,8$ remainder bits;
no new financial or smoothing assumption follows from this extension.
The fresh scalar tables still require checked rounding, monotonicity,
input domains and whole-graph absence of modular overflow.

\begin{lemma}[Literal selector count]
\label{lem:raw-selector-count}
Suppose a unary lookup has $b_e+1$ address bits and exactly the first
$N+1$ words are nonzero, where $N=2^{b_e}$ and $b_e\geq1$.
The retained recursively pruned selector tree visits exactly $N+b_e$
internal nodes. Its logical T count is
\[
 \tau_{\rm look}(b_e)=
 \begin{cases}
 14(N+b_e),&\text{exact seven-T Toffoli architecture},\\
 8(N+b_e),&\text{unitary eight-T selector-pair architecture}.
 \end{cases}
\]
\end{lemma}
\begin{proof}
The root partitions the $2N$ address values into two halves. Its left
subtree contains all $N$ nonzero words and is a full binary tree with
$N-1$ internal nodes. The right subtree contains the single nonzero word
at address $N$ and consequently visits a path with $b_e$ internal nodes.
The root itself contributes one, giving $1+(N-1)+b_e=N+b_e$.
Each visited node has exactly one selector compute and one selector
erase. In the first architecture these are two seven-T Toffolis. In the
second they are a four-T clean-target compute and its four-T unitary
inverse. All changes of branch, word writes and root operations are
Clifford gates. The stated counts therefore follow from the literal
retained sequence, with no measurement or free lookup assumption.
\end{proof}

All packed narrow-exponential words are nonzero because their base
ordinates are at least the exact integer representing one. The final
endpoint word is also present; it has zero interpolation delta and its
nonzero base. Thus Lemma~\ref{lem:raw-selector-count} applies.
Each exponential node contains two lookup invocations, as its private
arithmetic is cleaned. The two payoff nodes and the subsequent cleanup
of the complete preparation give eight such invocations in $A$.
Let $a^{\rm fixed}_{j,k}$ be the retained logical T count of all other
blocks and the comparator in case $j$, normal width $k$ and the stated
architecture, including their cleanup. Their arguments, word widths,
tables and emitted sequences are independent of $b_e$. Discarding every
other T gate inside the exponential nodes gives
\[
 T_A(b_e)\geq a^{\rm fixed}_{j,k}+c(2^{b_e}+b_e),
 \qquad c=112\text{ or }64.
\]
This is valid even when a wider address shortens its interpolation
multiplier, because that multiplier has been discarded completely from
the lower bound. It relies on this literal graph, without a hypothetical
optimisation which shares the two scalar evaluations.

Let $A_j=a_{+,j}$ be the unchanged outer mass upper bound, $e_{0,j}$
the coefficient radius, $\beta_j$ the financial-transfer upper bound,
$a_{r,j}$ the outer mass radius and $[l_j,u_j]$ the retained ideal
finite-mean interval. Give a competing wider implementation the
optimistic error $\eta=0$, and define
\[
 d_{0,j}=e_{0,j}+\beta_j+a_{r,j},\qquad
 \varepsilon_{0,j}=\frac{1/100-d_{0,j}}{A_j}.
\]
Every admitted implementation has deterministic error at least $d_{0,j}$
and hence probability tolerance at most $\varepsilon_{0,j}$. Its enlarged
probability upper prior is at least $u_j$. The controller's variance
upper rule
\[
 v(u)=\begin{cases}u(1-u),&0\leq u\leq1/2,\\1/4,&1/2<u\leq1\end{cases}
\]
is nondecreasing, as is its outward dyadic square-root rule. Therefore
let $M_{0,j}$ be the first supported power of two satisfying the retained
canonical-QPE sufficient amplitude bound at $(\varepsilon_{0,j},u_j)$.
Every active admitted plan has $M\geq M_{0,j}$. The retained inequalities
$A_j(u_j-l_j)/2>1/100-d_{0,j}$ and $u_j>2\varepsilon_{0,j}$ exclude
both controller forms of a zero-shot prior decision. These are checked
exactly for each of the four cases.

For confidence amplification, set $\pi_+=355/113$ and give the inverse
Fourier transform zero operator error. The controller's maximal certified
single-shot success probability is then $s_0=8/\pi_+^2$. Exact binomial
calculation shows that the first odd repeat count with majority failure
at most $1/3600$ is $r_0=25$. For a smaller success probability the
majority-failure probability is nonincreasing in the success parameter
and is therefore no better. This monotonicity follows, for example, by
coupling each success indicator as $\mathbf1\{U_i\leq s\}$ with common
independent uniform $U_i$. Every active plan thus has at least $r_0$
repetitions. This is a constraint of this certification controller;
it is not a lower bound on all conceivable amplitude estimators.

One canonical-QPE repetition calls $A$ once initially and calls both
$A$ and $A^\dagger$ in each of its $M-1$ controlled Grover steps.
Ignoring all reflection and Fourier T gates gives the complete bound
\[
 T_{\rm complete}(b_e)
 \geq r_0(2M_{0,j}-1)
       \left(\min_{k\in\{8,10,12\}}a^{\rm fixed}_{j,k}
             +c(2^{b_e}+b_e)\right).
\]
This lower bound concerns the fully specified logical resource policy;
it does not convert sufficient monetary-error bounds into an algorithmic
quantum query lower bound.

The exact retained fields give $M_{0,j}=256,512,512,512$ in case order.
For the less expensive unitary architecture the fixed-block minima are
$1\,494\,706$, $3\,129\,342$, $3\,348\,190$ and $7\,808\,330$ T gates,
all attained at normal width eight. At address width fifteen the lower
complete counts are respectively
\[
45\,898\,249\,950,\quad133\,692\,136\,050,\quad
139\,289\,173\,650,\quad253\,357\,254\,150.
\]
Each already exceeds an admitted earlier raw-payoff cap in its own case.
The exact seven-T architecture has still larger corresponding bounds.
Since $2^{b_e}+b_e$ is strictly increasing, all feasible widths fifteen
through twenty-two are excluded from improving those incumbents under
the stated fixed choices. Widths beyond twenty-two are outside this
fixed-point interpolator's input-resolution domain.

This exclusion does not cover omitted lower address widths
$\{1,2,3,5,7,9\}$, unequal widths for the two payoff nodes, other normal
widths, shared arithmetic, different word precisions, other lookup
architectures or other mean estimators. It can resolve an upper-boundary
concern for the registered extension without supporting a claim of a
global precision optimum, generic quantum superiority, theoretical
priority or practical quantum advantage.
